% chapter01.tex - 引言
\chapter{引言 (Introduction)}

\section{微分几何与物理学的关系}

微分几何是数学中研究光滑流形上的几何结构的分支，它为现代物理学提供了不可或缺的语言和工具。从牛顿力学到爱因斯坦的广义相对论，从麦克斯韦电磁理论到杨--米尔斯规范场论，物理学家发现了一个深层的真理：\textbf{自然界的定律本质上是几何的}。

\begin{note}[物理动机]
为什么物理学需要微分几何？答案在于一个看似平凡却极其深刻的原理：物理定律应当与坐标系的选择无关——这就是\textbf{广义协变性原理}。当我们选择不同的坐标系来描述同一个物理系统时，物理定律的形式必须保持不变。微分几何恰好提供了描述这种与坐标无关的几何结构的完整数学工具。
\end{note}

要理解这一思想的核心，我们可以从一个简单的问题出发：什么是"直线"？在欧几里得空间中，我们自然地回答"两点之间最短的路径"。但在弯曲的曲面上——比如球面——这个简单的概念需要重新审视。球面上两点之间的最短路径是连接它们的大圆弧，而非嵌入空间中的"笔直"线段。这个简单的观察揭示了一个关键事实：\textbf{几何结构决定了物体如何在空间中运动}。

\begin{figure}[htbp]
\centering
\begin{tikzpicture}[scale=1.3]
    % 球面
    \shade[ball color=blue!15,opacity=0.6] (0,0) circle (1.8);
    \draw (0,0) circle (1.8);
    % 赤道大圆
    \draw[red,thick] (-1.8,0) arc (180:360:1.8 and 0.5);
    \draw[red,thick,dashed] (-1.8,0) arc (180:0:1.8 and 0.5);
    % 两个点
    \fill[red] (-1.2,0.7) circle (0.06) node[above left] {$p$};
    \fill[red] (1.2,-0.5) circle (0.06) node[below right] {$q$};
    % 大圆弧
    \draw[red,thick] (-1.2,0.7) .. controls (-0.5,1.3) and (0.5,-0.1) .. (1.2,-0.5);
    \node[red] at (0.1,0.9) {测地线（大圆）};
    % 标注
    \node at (-2.0,-1.3) {球面 $S^2$};
\end{tikzpicture}
\caption{球面上两点 $p$ 和 $q$ 之间的最短路径是连接它们的大圆弧（测地线），而非三维嵌入空间中的"直线"。在弯曲空间中，几何结构决定了什么是"直的"。}
\label{fig:sphere_geodesic}
\end{figure}

在牛顿力学中，引力被看作一种"力"，它以超距作用的方式吸引物体。但这一描述有一个根本性的问题：为什么\textbf{引力质量} $m_G$（出现在万有引力定律 $F = G M m_G / r^2$ 中）恰好等于\textbf{惯性质量} $m_I$（出现在牛顿第二定律 $F = m_I a$ 中）？牛顿力学对此没有解释，它只能将此视为一个巧合。事实上，现代 E\"ot-Wash 实验已经以 $10^{-15}$ 的相对精度验证了 $m_G = m_I$ —— 这是物理学中最精确验证的等价性之一。

\medskip

\textbf{从 $m_G=m_I$ 到时空弯曲的逻辑链。} 这一等价性的深远意义可以通过以下推理来理解：
\begin{enumerate}
    \item 由于 $m_G=m_I$，所有物体在引力场中以相同的加速度下落——无论其质量或组成如何。
    \item 因此，一个自由下落的观察者无法通过任何局部实验来探测引力场——引力在局部上可以被"变换掉"。
    \item 这意味着引力不是一种普通的力，而是时空本身的结构属性。任何试图将引力纳入平坦时空的努力都会与等效原理相矛盾。
    \item 时空必须被赋予弯曲的几何结构，使得自由粒子沿测地线运动，而这一运动的表观加速度恰好解释了"引力"现象。
\end{enumerate}
这一逻辑链——从 $m_G=m_I$ 出发，经过局部不可分辨性，最终导出时空弯曲——是广义相对论最深刻的思想主线，也是为什么微分几何作为描述弯曲空间的数学语言成为理论物理学的核心工具。

\medskip

爱因斯坦给出了这一"巧合"的深层次解释：\textbf{引力不是一种力，而是时空弯曲的表现}。在弯曲的时空中，自由下落的物体沿测地线——即弯曲时空中对应的"直线"——运动。这一革命性的观念构成了广义相对论的基础，而描述弯曲时空的数学语言正是微分几何。

\medskip

这里值得一提的是\textbf{马赫原理 (Mach's Principle)}。马赫认为，物体的惯性并非其内禀属性，而是由宇宙中所有其他物质共同决定的——一个孤立粒子在空无一物的宇宙中将没有惯性。虽然爱因斯坦深受这一思想的启发，并希望在广义相对论中实现它，但最终的场方程并未完全体现马赫原理（例如，存在非马赫型的德西特宇宙解）。然而，广义相对论确实部分体现了马赫的精神：物质的分布（通过 $T_{\mu\nu}$）决定了时空的几何（通过 $g_{\mu\nu}$），而时空几何又决定了物体的惯性运动。这一物质-几何的相互决定关系成为相对论宇宙学的核心主题，至今仍然活跃在量子引力研究的前沿。

\begin{note}[物理意义：引力的几何本质]
爱因斯坦的洞见可以总结为"引力即几何"。这改变了我们对宇宙的根本理解：
\begin{itemize}
    \item 引力场 = 时空曲率（度规张量 $g_{\mu\nu}$ 描述）
    \item 自由落体 = 沿测地线运动（弯曲时空的"直线"）
    \item 潮汐力 = 相邻测地线的相对加速度（由曲率张量描述）
\end{itemize}
这一框架将引力从一种神秘的"力"转化为可精确计算的几何效应。
\end{note}

\section{历史发展}

微分几何与广义相对论的发展史是一个数学与物理学相互促进的经典案例。

\begin{figure}[htbp]
\centering
\begin{tikzpicture}[
    scale=0.85,
    timeline/.style={->,thick,blue!60},
    event/.style={fill=blue!50,text=white,rounded corners=3pt,inner sep=3pt,font=\small,align=center},
    math_event/.style={fill=green!50!black,text=white,rounded corners=3pt,inner sep=3pt,font=\small,align=center},
    phys_event/.style={fill=red!50,text=white,rounded corners=3pt,inner sep=3pt,font=\small,align=center}
]
    % Time axis
    \draw[timeline] (0,0) -- (16,0) node[right] {时间};
    
    % Gauss 1827
    \draw (0.5,0.3) -- (0.5,-1.5);
    \node[event] at (0.5,-1.8) {Gauss 1827\\Theorema Egregium};
    \node at (0.5,0.6) {1827};
    
    % Riemann 1854
    \draw (3,0.3) -- (3,-1.5);
    \node[event] at (3,-1.8) {Riemann 1854\\$n$维弯曲空间};
    \node at (3,0.6) {1854};
    
    % Ricci & Levi-Civita 1900
    \draw (8,0.3) -- (8,-1.5);
    \node[event] at (8,-1.8) {Ricci \& Levi-Civita 1900\\张量分析（绝对微分学）};
    \node at (8,0.6) {1900};
    
    % Grossmann + Einstein 1913
    \draw (12.5,0.3) -- (12.5,-1.5);
    \node[event] at (12.5,-1.8) {Einstein \& Grossmann\\1913 黎曼几何+物理};
    \node at (12.5,0.6) {1913};
    
    % Einstein 1915
    \draw (15,1.8) -- (15,-0.2);
    \node[event] at (15,2.2) {Einstein 1915\\$G_{\mu\nu}=8\pi G T_{\mu\nu}$};
    \node at (15,1.7) {1915};
    
    % Schwarzschild 1916
    \draw (15,-0.8) -- (15,-2.2);
    \node[event] at (15,-2.5) {Schwarzschild 1916\\第一个黑洞解};
    
    % Weyl, Cartan
    \draw (13.5,-3.2) node[event] {Weyl \& Cartan\\纤维丛与联络};
    
    % Wheeler 1967
    \draw (10,-3.2) node[event] {Wheeler 1967\\"黑洞"命名};
    
    % Penrose & Hawking
    \draw (5.5,-3.2) node[event] {Penrose \& Hawking\\1965-70 奇点定理};
    
    % LIGO 2015
    \draw (1,-3.2) node[event] {LIGO 2015\\引力波探测};
    
    % Labels
    \node[font=\footnotesize\itshape,gray] at (15.5,-3.7) {... 至今};
\end{tikzpicture}
\caption{微分几何与广义相对论的发展简史。绿色代表数学进展，红色代表物理进展。两者相互促进：黎曼几何为广义相对论提供了数学工具，而广义相对论反过来又刺激了现代微分几何的发展。}
\label{fig:history_timeline}
\end{figure}

下面我们对各个关键里程碑进行更详细的叙述。

\begin{enumerate}
    \item \textbf{高斯 (Gauss, 1827) —— 曲面的内蕴几何}：高斯在研究二维曲面时提出了一个划时代的观点：曲面的几何性质可以由曲面自身上的测量来决定，而不需要参考其嵌入的高维空间。他在《关于曲面的一般研究》中证明的 \textbf{Theorema Egregium}（绝妙定理）表明，高斯曲率 $K$ 完全由曲面上的度规系数 $g_{ij}$ 及其偏导数决定，是一个内蕴量。这意味着生活在曲面上的"二维生物"可以通过测量三角形的内角和来探测曲面的弯曲，而无需知晓第三维的存在——这蕴含了现代物理学中"生活在时空中而不需要外部参考"的思想 \cite{Carroll1997}。
    
    \item \textbf{黎曼 (Riemann, 1854) —— 高维黎曼几何}：黎曼将高斯的思想推广到任意维数。在他著名的就职演讲《论作为几何学基础的假设》中，黎曼定义了 $n$ 维弯曲空间的度规：
    \begin{equation}
        \dd s^2 = g_{\mu\nu}(x) \dd x^\mu \dd x^\nu,
    \end{equation}
    并提出曲率完全由度规张量 $g_{\mu\nu}$ 及其导数决定。这一工作为后来爱因斯坦的广义相对论提供了必需的数学框架。
    
    \item \textbf{黎曼曲率张量的发现}：黎曼定义了四阶曲率张量 $R^\lambda{}_{\mu\nu\rho}$，其物理意义是：当一个向量沿无限小闭合回路平行移动一周后，它不回初始方向的偏离量。这正是引力的几何诠释的核心——\textbf{曲率就是引力场}。曲率非零意味着存在潮汐力。
    
    \item \textbf{里奇 (Ricci-Curbastro) 与列维--奇维塔 (Levi-Civita, 1900) —— 张量分析}：里奇和列维--奇维塔发展了所谓的"绝对微分学"（即今天所称的张量分析），将黎曼几何的数学语言系统化为便于物理应用的形式。列维--奇维塔引入了\textbf{平行移动}（parallel transport）的概念，为弯曲空间中的"方向比较"提供了数学基础。他还提出引力的标量理论候选者应该是黎曼标量曲率 $R$——这一步虽然不完整，但为后来爱因斯坦的工作指明了方向 \cite{Wald1984}。
    
    \item \textbf{爱因斯坦与格罗斯曼 (Einstein \& Grossmann, 1912--1915) —— 寻找场方程}：1912 年，爱因斯坦从布拉格返回苏黎世，找到了他的老友马塞尔·格罗斯曼——苏黎世联邦理工学院的数学教授。据说爱因斯坦对他说的第一句话是："格罗斯曼，你必须帮助我，否则我会发疯！" \cite{Wald2019}。格罗斯曼查阅了图书馆的数学文献，发现了黎曼、里奇和列维--奇维塔的工作，意识到这正是描述弯曲时空所需的数学工具。1913 年，两人合作发表了《广义相对论和引力理论纲要》，虽然其中的场方程还不完整，但它已经将物理（等效原理和广义协变性）与数学（黎曼几何）紧密地联系在了一起 \cite{Tong_GR_LectureNotes}。
    
    \item \textbf{爱因斯坦 (1915) —— 场方程的最终形式}：经过两年的艰苦努力，爱因斯坦于 1915 年 11 月在普鲁士科学院连续四周提交了四篇论文，最终完成了场方程的正确形式：
    \begin{equation}
        R_{\mu\nu} - \frac{1}{2}R g_{\mu\nu} = \frac{8\pi G}{c^4} T_{\mu\nu},
    \end{equation}
    并在此框架下解释了水星近日点的反常进动（每世纪 43 角秒，此前是经典力学无法解释的谜题），成功地预言了光线在太阳引力场中的偏折（1919 年日食观测证实，使爱因斯坦一夜之间成为世界名人）。爱因斯坦场方程简洁而深刻：\textbf{左侧是时空几何（爱因斯坦张量），右侧是物质分布（能量-动量张量）}。
    
    \item \textbf{施瓦西 (Schwarzschild, 1916) —— 第一个精确解}：在爱因斯坦发表场方程后不到两个月，德国天文学家施瓦西——当时正在第一次世界大战的东线服役——找到了球对称真空解的精确表达式，即施瓦西度规。这个解如今被理解为一个非旋转黑洞的时空几何，开启了黑洞物理的研究。施瓦西在论文发表后不久就因病去世，年仅 43 岁。
    
    \item \textbf{外尔与嘉当 —— 纤维丛与联络}：外尔 (Hermann Weyl) 和嘉当 (Élie Cartan) 在 20 世纪初进一步推广了联络的概念，将黎曼几何从度量联络扩展到了一般的仿射联络，为现代规范场论奠定了数学基础。嘉当还发展了\textbf{活动标架法} (moving frame method) 和外微分计算，使得曲率的计算更加系统和高效。有趣的是，外尔的标度不变几何虽然作为引力理论失败了，但后来作为电磁规范理论的先驱而获得了新生。
    
    \item \textbf{惠勒 (Wheeler, 1955--1970) —— "几何动力学"与黑洞命名}：约翰·惠勒重新点燃了物理学界对广义相对论的兴趣。他在 1955 年为这门学科命名为"几何动力学" (geometrodynamics)，强调其几何本质。1967 年，惠勒在一次讲座中首次使用了"黑洞" (black hole) 这一术语，使其从一个物理学概念转变为流行文化的一部分。
    
    \item \textbf{彭罗斯与霍金 (1965--1970) —— 奇点定理}：彭罗斯 (1965) 和霍金 (1966--1970) 证明了奇点定理——在非常一般的条件下，引力坍缩必然导致时空奇点的形成。彭罗斯因此获得 2020 年诺贝尔物理学奖 \cite{Penrose1965}。这些定理标志着广义相对论自身包含了对经典框架彻底失效的预测，从而表明了量子引力修正的必要性。
    
    \item \textbf{LIGO (2015) —— 引力波的直接探测}：2015 年 9 月 14 日，LIGO 首次直接探测到来自双黑洞并合的引力波。这一里程碑式的发现 \cite{Abbott2016} 不仅验证了爱因斯坦的最后一个未经验证的主要预言，还开启了"引力波天文学"的新纪元。2017 年诺贝尔物理学奖授予了 LIGO 的三位奠基人。
\end{enumerate}

\subsection{数学与物理的相互滋养}

微分几何与广义相对论的关系史是数学和物理学相互促进的典范，堪称科学史上的"正反馈循环"。这一循环大致可分为三个阶段：

\medskip

\textbf{第一阶段：数学先行（1827--1915）。} 高斯的内蕴几何理论和黎曼的$n$维弯曲空间理论完全是在纯数学的驱动下发展的。黎曼在1854年的就职演讲中提出的度规概念和曲率张量，其动机来自对几何学基础的哲学思考，而非任何物理应用。事实上，当时没有任何物理理论需要用到这些工具——直到爱因斯坦登场。

\medskip

\textbf{第二阶段：物理整合（1912--1915）。} 当爱因斯坦为广义相对论寻找数学语言时，格罗斯曼引导他发现了黎曼、里奇和列维--奇维塔的工作。爱因斯坦的物理洞察力与这套现成的数学工具产生了惊人的化学反应——场方程本身就是物理学与数学完美结合的产物。有趣的是，爱因斯坦起初对张量分析持怀疑态度，认为它是"多余的学问"（\textit{überflüssige Gelehrsamkeit}），但最终他不得不承认，没有这套数学工具，广义相对论根本不可能被写下来。

\medskip

\textbf{第三阶段：物理反哺数学（1916--至今）。} 广义相对论的成功极大地刺激了微分几何本身的发展。以下是几个典型的例子：

\begin{itemize}
    \item \textbf{全局微分几何的兴起}：广义相对论提出了经典黎曼几何未曾面对的问题——时空的全局结构是什么？奇点是否不可避免？这促使数学家发展了一套全局微分几何理论。彭罗斯 (Roger Penrose) 在1965年引入的\textbf{旋量方法 (spinor methods)}和共形紧化技术，不仅证明了第一个奇点定理，还为现代微分几何开辟了全新的方向。Geroch、Hawking 和 Ellis 等人将\textbf{拓扑学}系统地引入广义相对论，用因果结构和整体双曲性等概念来刻画时空的全局性质。
    
    \item \textbf{爱因斯坦流形与几何分析}：场方程 $R_{\mu\nu} = \Lambda g_{\mu\nu}$（真空+宇宙常数）在数学中定义了\textbf{爱因斯坦流形 (Einstein manifold)}。对这类流形的研究成为现代微分几何的重要分支。丘成桐 (Shing-Tung Yau) 在证明 Calabi 猜想（1976）时发展的偏微分方程技术，后来成为几何分析的核心工具，并为他赢得了菲尔兹奖。
    
    \item \textbf{正质量定理}：丘成桐 (1979) 和 Witten (1981) 分别用几何分析和旋量方法证明了广义相对论中的一个基本物理命题——在适当的能量条件下，孤立系统的总质量（ADM质量）非负。这一"物理定理"不仅具有深刻的物理意义，其证明过程中的数学技巧也推动了微分几何的发展。
    
    \item \textbf{规范场论与纤维丛}：广义相对论（引力规范理论）启发了杨--米尔斯规范场论（1954）的建立。反过来，规范场论的数学结构——纤维丛上的联络与曲率——又被数学家发展为现代微分几何的核心框架，并反过来深化了我们对引力理论本身的理解（例如，Ashtekar 变量和圈量子引力）。
    
    \item \textbf{引力波的数学理论}：引力波的数学描述——包括 Bondi-Sachs 能-动量的定义、渐近平坦时空的刻画、以及 Christodoulou 和 Klainerman 对爱因斯坦方程非线性稳定性的证明——都需要极其精深的微分几何和偏微分方程技术，这些工作在数学本身中也引起了广泛关注。
\end{itemize}

这一相互滋养的循环至今仍在继续：弦理论和圈量子引力对量子时空几何的需求正在推动新的数学发展，而数学中关于 Ricci 流、几何化猜想等领域的最新进展又为引力理论提供了新的视角。正如著名数学家 Michael Atiyah 所说："广义相对论是赋予微分几何以实质的物理理论——没有它，微分几何可能仍然只是数学家的智力游戏。"

\section{等效原理：引力与加速度的等价性}
\label{sec:equivalence_principle}

爱因斯坦将他在 1907 年做出的关键洞察称为自己"一生中最快乐的思想"（\textit{der glücklichste Gedanke meines Lebens}）：\textbf{在自由落体的参考系中，引力场局部消失}。这一思想可以用著名的电梯实验来说明。

\begin{figure}[htbp]
\centering
\begin{tikzpicture}[scale=1.0]
    % === Left: 地球上的静止电梯 ===
    \draw[thick] (0,0) rectangle (2.5,3.5);
    \node at (1.25,3.7) {地球上};
    
    % 人
    \draw[thick] (1.0,0.3) circle (0.2); % head
    \draw[thick] (1.0,0.1) -- (1.0,-0.6); % body
    \draw[thick] (1.0,-0.3) -- (0.7,-0.7); % left leg
    \draw[thick] (1.0,-0.3) -- (1.3,-0.7); % right leg
    
    % 重力箭头
    \draw[->,blue,very thick] (1.5,1.5) -- (1.5,0.3) node[midway,right,font=\small] {重力};
    \draw[->,red,very thick] (0.7,-0.7) -- (0.7,-1.8) node[midway,right,font=\small] {$mg$};
    \draw[->,blue!50,very thick] (1.25,-0.7) -- (1.25,0.2) node[midway,left,font=\small] {支持力};
    
    % 地球
    \draw[thick,brown!70] (-1,-2.5) arc (180:360:3.5 and 1.0);
    \draw[thick,brown!70] (-1.5,-2.5) -- (5.5,-2.5);
    \node at (4.5,-2.8) {地球};
    
    % === Center: 太空中的加速电梯 ===
    \draw[thick] (5.5,0) rectangle (8,3.5);
    \node at (6.75,3.7) {太空中};
    
    % 人
    \draw[thick] (6.5,0.3) circle (0.2);
    \draw[thick] (6.5,0.1) -- (6.5,-0.6);
    \draw[thick] (6.5,-0.3) -- (6.2,-0.7);
    \draw[thick] (6.5,-0.3) -- (6.8,-0.7);
    
    % 加速箭头
    \draw[->,red,very thick] (7.0,1.5) -- (7.0,0.3) node[midway,right,font=\small] {惯性力};
    \draw[->,blue!50,very thick] (6.0,2.0) -- (6.0,3.2) node[midway,left,font=\small] {$\mathbf{a}=g$};
    \node at (6.75,1.0) {加速上升};
    
    % Stars
    \foreach \x/\y in {5.6/3.0,7.8/2.8,6.0/2.5} {
        \fill[gray] (\x,\y) circle (0.03);
    }
    
    % === Right: 自由下落电梯 ===
    \draw[thick] (10.5,2.0) rectangle (13.0,5.5);
    \node at (11.75,5.7) {自由下落};
    
    % 漂浮的人
    \draw[thick] (11.5,3.5) circle (0.2);
    \draw[thick] (11.5,3.3) -- (11.5,2.6);
    \draw[thick] (11.5,2.9) -- (11.2,2.5);
    \draw[thick] (11.5,2.9) -- (11.8,2.5);
    
    % 下落箭头
    \draw[->,red,very thick] (12.0,5.2) -- (12.0,3.0) node[midway,right,font=\small] {$g$};
    \node at (11.7,4.0) {感受不到重力！};
    
    % Earth
    \draw[thick,brown!70] (9,-0.5) arc (180:360:3.5 and 1.0);
    \draw[thick,brown!70] (8.5,-0.5) -- (13.5,-0.5);
    
    % Labels
    \node at (1.25,-3.5) {(a) 地球表面};
    \node at (6.75,-3.5) {(b) 太空加速};
    \node at (11.75,-3.5) {(c) 自由下落};
\end{tikzpicture}
\caption{等效原理的电梯思想实验。(a) 地球上静止的电梯中的人感受到重力。(b) 太空中以加速度 $g$ 加速上升的电梯——内部的人感受到完全相同的"重力"。(c) 在自由下落的电梯中，引力的效应被完全消除——人漂浮起来，感受不到"重力"。由此，在局部无法区分引力场与加速参考系。}
\label{fig:elevator_experiment}
\end{figure}

\subsection{等效原理与引力时间膨胀的定量推导}

等效原理的一个最直接、最可观测的推论是\textbf{引力时间膨胀}（也称为引力红移）。下面我们通过一个定量的推导来展示等效原理如何预测这一效应。这是爱因斯坦在1907年做出的关键计算之一，也是等效原理威力的绝佳例证。

\medskip

\noindent\textbf{步骤一：加速电梯中的多普勒频移。} 考虑一台在无引力太空中以恒定加速度 $\mathbf{a}$ 向上加速的电梯。电梯的高度为 $h$，底部（floor）位于 $z=0$，顶部（ceiling）位于 $z=h$。在 $t=0$ 时刻，底部的光源向顶部发出一个频率为 $\nu_0$ 的光信号。

在光信号从底部传播到顶部的时间 $\Delta t \approx h/c$ 内（假设 $h$ 很小，加速度效应在一阶近似下处理），电梯顶部的接收器由于加速运动获得了一个额外的速度：
\begin{equation}
    \Delta v = a \cdot \Delta t \approx a \cdot \frac{h}{c}.
\end{equation}
根据经典多普勒效应公式（非相对论极限下，$v \ll c$），接收器测量到的频率为：
\begin{equation}
    \nu_{\text{obs}} = \nu_0 \left(1 - \frac{\Delta v}{c}\right) = \nu_0 \left(1 - \frac{a h}{c^2}\right).
\end{equation}
因此，在加速电梯中，顶部接收到的光发生了\textbf{红移}（频率降低）：
\begin{equation}
    \frac{\Delta \nu}{\nu_0} \equiv \frac{\nu_{\text{obs}} - \nu_0}{\nu_0} = -\frac{a h}{c^2}.
\end{equation}

\medskip

\noindent\textbf{步骤二：通过等效原理转换为引力场。} 现在应用爱因斯坦等效原理：一个以加速度 $a$ 向上加速的电梯，在局部上与一个处于均匀引力场 $g = a$ 中静止的电梯不可区分。将加速电梯中的结果翻译为引力场中的物理：

\begin{itemize}
    \item 在加速电梯中，顶部接收到的光频率低于底部发出的频率——因为顶部在"逃离"光源。
    \item 在引力场中，这意味着：\textbf{在引力场中不同高度的时钟以不同的速率运行。位于引力势较低处（即更靠近引力源，引力更强处）的时钟走得较慢。}
\end{itemize}

将 $a$ 替换为引力加速度 $g$，并注意到 $gh = \Delta \Phi$（引力势差），我们得到：
\begin{equation}
    \boxed{\frac{\Delta \nu}{\nu_0} = -\frac{g h}{c^2} = -\frac{\Delta \Phi}{c^2}}.
\end{equation}
这是引力红移的基本公式：当光从引力势较低处（$\Phi$ 更负）向较高处传播时，频率降低——即发生红移。

\medskip

\noindent\textbf{步骤三：用度规语言表述。} 在广义相对论的完整框架中，这一结果可以用度规张量来精确表达。考虑静态弱场度规：
\begin{equation}
    \dd s^2 = -\left(1 + \frac{2\Phi(\mathbf{x})}{c^2}\right) c^2 \dd t^2 + \left(1 - \frac{2\Phi(\mathbf{x})}{c^2}\right) \dd \mathbf{x}^2,
\end{equation}
其中 $\Phi(\mathbf{x})$ 是牛顿引力势（对于地球表面附近，$\Phi = gz$）。对于两个静止在不同高度 $z_1$ 和 $z_2$ 的时钟，其固有时 $d\tau = \sqrt{-g_{00}}\,dt$ 的比率为：
\begin{equation}
    \frac{\dd \tau_2}{\dd \tau_1} = \sqrt{\frac{g_{00}(z_2)}{g_{00}(z_1)}} = \sqrt{\frac{1 + 2\Phi_2/c^2}{1 + 2\Phi_1/c^2}} \approx 1 + \frac{\Phi_2 - \Phi_1}{c^2} = 1 + \frac{g(z_2 - z_1)}{c^2}.
\end{equation}
对于 $z_2 > z_1$（高处），$\dd\tau_2 / \dd\tau_1 > 1$——高处的时钟走得\textbf{更快}。

\medskip

\noindent\textbf{实验验证。} 这一效应已被多个实验精确验证：
\begin{itemize}
    \item \textbf{Pound--Rebka 实验 (1959)}：利用 M\"ossbauer 效应在哈佛大学 Jefferson 实验室的塔楼（高度差仅 $22.5\,\text{m}$）中测量了地球引力场中的引力红移，精度约为 $1\%$。预测的红移相对大小为 $\Delta \nu / \nu \approx 2.5 \times 10^{-15}$，极其微小，但实验成功探测到了这一效应。
    \item \textbf{Gravity Probe A (1976)}：将氢原子钟搭载在火箭上送到 $10{,}000\,\text{km}$ 的高度，以约 $7\times 10^{-5}$ 的精度验证了引力红移公式。
    \item \textbf{GPS 系统（全球定位系统）}：GPS 卫星在约 $20{,}200\,\text{km}$ 的高度运行。由于引力时间膨胀，卫星上的时钟每天比地面时钟快约 $45\,\mu\text{s}$（此外还有狭义相对论速度时间膨胀导致的每天慢约 $7\,\mu\text{s}$，净效应为每天快 $38\,\mu\text{s}$）。如果不对此进行修正，GPS 的定位误差将以每天约 $11\,\text{km}$ 的速度累积——这是广义相对论在日常生活技术中的直接应用。具体计算见第 \ref{sec:examples} 节的算例 \ref{para:gps}。
\end{itemize}

这个推导展示了等效原理的强大力量：仅从加速度和经典多普勒效应的考虑出发，我们就定量地预测了引力时间膨胀——一个纯粹的广义相对论效应——而不需要用到任何弯曲时空的数学！这正是爱因斯坦将其称为"一生中最快乐的思想"的原因。

\begin{definition}[弱等效原理 (WEP)]
在引力场中，所有物体——无论其质量或内部组成——都以相同的加速度下落。用牛顿力学的语言：惯性质量 $m_I$ 严格等于引力质量 $m_G$。这一定律最早由伽利略通过比萨斜塔思想实验提出，后由 E\"otv\"os (1889)、Dicke (1964) 和现代 E\"ot-Wash 实验以越来越高精度验证。
\end{definition}

\begin{definition}[爱因斯坦等效原理 (EEP)]
在足够小的时空区域内，物理定律的形式与其在无引力场的惯性参考系中的形式完全相同。也就是说：\textbf{在任意引力场中的每一点，都存在一个局部惯性系，使得狭义相对论的物理定律在此参考系中成立。}
\end{definition}

爱因斯坦等效原理有几个核心推论：
\begin{itemize}
    \item \textbf{引力红移}：在引力场中不同位置，时钟运行的速率不同。位于引力势较低（引力更强）处的时钟走得较慢。
    \item \textbf{光线偏折}：光线在引力场中弯曲传播。因为光在局部惯性系中沿直线传播，但在全局弯曲时空中，这些"局部直线"拼接起来形成了一条弯曲的轨迹。
    \item \textbf{时空必须是弯曲的}：引力不能在平坦的闵可夫斯基时空中被一致地描述。等效原理要求全局时空的几何结构是弯曲的——这正是黎曼几何的用武之地。
\end{itemize}

\begin{figure}[htbp]
\centering
\begin{tikzpicture}[scale=1.1]
    % Flat grid background (Minkowski)
    \draw[gray!30,step=0.5] (-3,-2.5) grid (3,2.5);
    \draw[gray!50] (-3,0) -- (3,0);
    \draw[gray!50] (0,-2.5) -- (0,2.5);
    
    % Matter blob (Earth)
    \fill[blue!30] (0,0) ellipse (0.6 and 0.4);
    \node[blue!70] at (0,0.7) {物质};
    
    % Curved grid (spacetime) - using manual coordinates to avoid foreach/plot conflicts
    \draw[red!40,smooth] (-2.6,0.5) .. controls (-1.5,0.3) .. (0,-0.1) .. controls (1.5,0.3) .. (2.6,0.5);
    \draw[red!40,smooth] (-2.7,-0.2) .. controls (-1.5,-0.1) .. (0,-0.5) .. controls (1.5,-0.1) .. (2.7,-0.2);
    \draw[red!40,smooth] (-2.8,-0.8) .. controls (-1.5,-0.5) .. (0,-0.7) .. controls (1.5,-0.5) .. (2.8,-0.8);
    \draw[red!40,smooth] (-2.9,-1.4) .. controls (-1.5,-0.9) .. (0,-0.9) .. controls (1.5,-0.9) .. (2.9,-1.4);
    \draw[red!40,smooth] (-2.5,1.1) .. controls (-1.5,0.7) .. (0,0.4) .. controls (1.5,0.7) .. (2.5,1.1);
    \draw[red!40,smooth] (-2.6,1.7) .. controls (-1.5,1.1) .. (0,0.8) .. controls (1.5,1.1) .. (2.6,1.7);
    \draw[red!40] (-2.5,-1.9) .. controls (-2.0,-1.6) and (2.0,-1.6) .. (2.5,-1.9);
    \draw[red!40] (-2.0,2.3) .. controls (-1.5,2.0) and (1.5,2.0) .. (2.0,2.3);
    \draw[red!40] (-1.5,-2.0) .. controls (-1.0,-1.8) and (1.0,-1.8) .. (1.5,-2.0);
    \draw[red!40] (-1.5,2.5) .. controls (-1.0,2.2) and (1.0,2.2) .. (1.5,2.5);
    
    % Curvature "dip"
    \draw[red,thick,domain=-2.5:2.5,samples=50,smooth]
        plot (\x, {-0.6*exp(-(\x)^2/1.5)});
    
    % Labels
    \node[blue!70] at (0,-2.2) {物质使时空弯曲};
    \node[red!70] at (3.5,1.8) {弯曲时空};
    
    % Test particle
    \draw[->,teal,thick] (2.2,1.0) .. controls (2.0,0.5) and (0.5,-0.2) .. (0.2,-0.55);
    \node[teal] at (3.0,0.8) {测地线};
\end{tikzpicture}
\caption{广义相对论的核心图景：物质（蓝色）导致时空弯曲（红色凹陷），而自由粒子沿着弯曲时空中的测地线运动（青色轨迹）。这就是惠勒总结的："时空告诉物质如何运动；物质告诉时空如何弯曲。"}
\label{fig:gr_curvature}
\end{figure}

\begin{note}[物理意义]
等效原理的深层意义在于：引力可以被"变换掉"——在局部自由下落参考系中，引力消失。这意味着引力不是一种普通的"力"，而是时空几何的一个方面。正如惠勒 (John Wheeler) 所总结的：\textbf{"时空告诉物质如何运动；物质告诉时空如何弯曲。"}
\end{note}

\section{广义协变性与广义相对论的基本原理}

\begin{property}[广义协变性原理]
物理定律在所有坐标系——不仅是惯性系——中都具有相同的形式。也就是说，物理定律的数学表达应该是张量方程：$T^{\mu_1\cdots}{}_{\nu_1\cdots} = 0$ 的形式在任何坐标变换下都不变。这确保了物理定律不依赖于描述它的具体方式。
\end{property}

\medskip

\noindent\textbf{广义协变性的双重含义。} 有必要区分"广义协变性"的两个不同层次：

\begin{itemize}
    \item \textbf{作为数学陈述的广义协变性：} 所有物理定律都可以写成张量方程的形式，因此在任意坐标变换下形式不变。这是一个关于理论\textit{表述方式}的断言——任何物理理论，原则上都可以通过引入适当的背景结构被写成广义协变的形式。例如，牛顿力学和狭义相对论都可以被重新表述为张量方程（只需引入平坦度规 $\eta_{\mu\nu}$ 作为背景场），尽管这样做可能并不自然。
    
    \item \textbf{作为物理原理的广义协变性：} 不存在\textit{先验}的、特殊的坐标系或参考系——时空本身没有"背景几何结构"，所有几何量（包括度规 $g_{\mu\nu}$）都是动力学的。这不是一个关于表述方式的断言，而是一个关于\textit{物理现实}的断言：时空的几何本身由物质分布决定，而不是预先给定的。
\end{itemize}

\medskip

\noindent\textbf{Kretschmann 的批评与爱因斯坦的回应。} 1917年，德国物理学家 Erich Kretschmann 提出了一个著名的批评：\textit{广义协变性作为一个选择物理理论的标准是空洞的}。他指出，\textbf{任何}物理理论——包括牛顿力学——都可以通过引入适当的变量被重写为广义协变的形式。因此，单独要求一个理论是"广义协变的"并不能将其与其他理论区分开来。

Kretschmann 的论点在数学上是正确的。但爱因斯坦的广义协变性远不止是一个形式要求。它的物理内容来自于与\textbf{等效原理}的结合：

\begin{center}
    \fbox{\parbox{0.9\textwidth}{
        \smallskip
        \textbf{爱因斯坦版本的广义协变性 = 广义协变性（数学要求）+ 等效原理（物理要求）}
        \smallskip
        
        等效原理告诉我们：在局部惯性系中，狭义相对论成立，且度规 $g_{\mu\nu}$ 约化为 $\eta_{\mu\nu}$。这一条件排除了那些仅仅通过"添加背景结构"而获得广义协变形式的理论——因为等效原理要求时空的局部几何必须是闵可夫斯基的，而这只有在度规本身是动力学的、能够响应物质分布的情况下才可能实现。
        \smallskip
    }}
\end{center}

换句话说，Kretschmann 的正确批评恰好凸显了为什么等效原理是广义相对论中不可或缺的第二个基石。单独一个广义协变性确实空洞——任何理论都可以被写成张量方程。但加上等效原理后，我们要求：\textit{在任意时空点的局部惯性系中，物理学回归狭义相对论，且这一局部惯性系由动力学的度规场 $g_{\mu\nu}$ 而非预先给定的背景结构来定义。} 正是这个组合使得广义相对论成为独一无二的引力理论。

\medskip

\noindent\textbf{现代视角：无背景结构 (Background Independence)。} 在当代量子引力的讨论中，广义协变性的物理核心被重新表述为\textbf{无背景结构原则 (background independence)}——一个基本的引力理论不应依赖于任何固定的、非动力学的时空背景。这一原则是圈量子引力、因果动力三角化等量子引力方案的核心指导原则，也是引力量子化之所以困难的根源——在量子场论中，我们习惯在有固定背景的时空中进行微扰计算，而广义相对论恰恰告诉我们，时空本身没有固定的背景。

广义相对论建立在两个基本假设之上：

\begin{enumerate}
    \item \textbf{等效原理}：在局部惯性系中，狭义相对论成立。这确保了引力在局部上可以被"加速参考系"取代。
    \item \textbf{广义协变性}：物理定律在所有坐标系中形式相同，因此可以用张量语言来表达。
\end{enumerate}

从这两个原理出发，我们可以构建广义相对论的完整框架：

\begin{itemize}
    \item 时空被建模为一个配备了洛伦兹号差度规 $g_{\mu\nu}$ 的四维光滑流形 $(\mathcal{M}, g)$。
    \item 度规 $g_{\mu\nu}$ 在局部惯性系中约化为闵可夫斯基度规 $\eta_{\mu\nu} = \operatorname{diag}(-1,+1,+1,+1)$。
    \item 自由粒子（不受非引力作用）沿测地线运动，测地线方程由度规的 Christoffel 符号决定。
    \item 物质的能量-动量张量 $T_{\mu\nu}$ 通过爱因斯坦场方程决定了时空的曲率，从而决定了度规 $g_{\mu\nu}$。
\end{itemize}

\begin{note}[物理意义：从"力"到"几何"的范式转变]
广义相对论带来的变革非同小可。在牛顿物理学中："一个物体受到引力，从而改变运动状态"；在广义相对论中："一个物体在弯曲时空中沿测地线运动，我们观察到它的轨迹弯曲，误以为有力作用在它身上"。换句话说，\textbf{引力不是一种相互作用，而是时空弯曲的几何效应}。
\end{note}

\section{本文的结构与阅读指南}

本文旨在从微分几何的数学基础出发，系统地引导读者理解广义相对论的核心思想和主要应用。本文假设读者已经掌握本科层次的微积分、线性代数、经典力学和电磁学，并具备基本的狭义相对论知识。

\begin{itemize}
    \item \textbf{第一部分：微分几何基础（第 \ref{chap:manifold}--\ref{chap:geodesic} 章）}
    \begin{itemize}
        \item 第 \ref{chap:manifold} 章建立流形的基本概念——将时空建模为光滑流形，是后面所有几何结构的基础
        \item 第 \ref{chap:tensor} 章引入张量及其代数运算——物理学定律的通用语言
        \item 第 \ref{chap:forms} 章讨论微分形式和外微分——弯曲时空上的积分与 Stokes 定理
        \item 第 \ref{chap:connection} 章引入联络与协变导数——如何在弯曲时空上"求导"
        \item 第 \ref{chap:curvature} 章讨论黎曼曲率张量——衡量时空弯曲程度的数学量
        \item 第 \ref{chap:geodesic} 章推导测地线方程——自由粒子如何在弯曲时空中运动
    \end{itemize}
    \item \textbf{第二部分：广义相对论（第 8--14 章）}
    \begin{itemize}
        \item 第 8 章回顾狭义相对论——广义相对论的平坦极限
        \item 第 9 章深入讨论等效原理与广义协变性——广义相对论的物理基石
        \item 第 10 章从作用量原理推导爱因斯坦场方程——时空几何如何与物质耦合
        \item 第 11 章给出施瓦西解与黑洞物理——广义相对论的第一个精确解
        \item 第 12 章讨论引力波——时空涟漪的物理
        \item 第 13 章介绍宇宙学基础——将广义相对论应用于整个宇宙
        \item 第 14 章简介数值相对论——用计算机求解爱因斯坦方程
    \end{itemize}
\end{itemize}

\begin{remark}
如果读者对抽象的数学定义感到不适，可以先阅读每个章节开头的"物理动机"部分，这些部分用直观的语言说明了为什么需要这些数学概念，然后再回到严格的定义。物理学家的方法从来都是：\textbf{先理解物理，再学数学。}正如 Dirac 所说："一条物理定律必须具备数学的美感。"——但美的理解始于直观。
\end{remark}

\section{算例：广义相对论的牛顿极限}
\label{sec:examples}

在深入微分几何之前，让我们通过一个具体算例来预览广义相对论与牛顿引力之间的联系。


\paragraph{算例 1.1：从测地线方程导出牛顿引力}
在弱场、慢速极限下，度规取 $g_{\mu\nu} = \eta_{\mu\nu} + h_{\mu\nu}$（$|h_{\mu\nu}| \ll 1$），其中 $h_{00} = -2\Phi$（$\Phi$ 是牛顿引力势），其余 $h_{ij} \approx 0$。类时测地线方程 $\dd^2 x^\mu/\dd\tau^2 + \Gamma^\mu{}_{\nu\rho} \dot x^\nu \dot x^\rho = 0$ 在低速近似 $\dot x^i \ll \dot x^0 \approx 1$ 下约化为：
\begin{equation}
    \frac{\dd^2 x^i}{\dd t^2} \approx -\Gamma^i{}_{00} = -\frac{1}{2}(\pd_i h_{00}) = -\pd_i \Phi.
\end{equation}
这就是 $\mathbf{a} = -\nabla \Phi$——牛顿第二定律！这个算例表明：\textbf{牛顿引力是广义相对论在弱场慢速极限下的近似}。

\begin{figure}[htbp]
\centering
\begin{tikzpicture}[scale=1.3]
    \shade[top color=blue!5,bottom color=blue!20,opacity=0.7] (0,0) ellipse (2.0 and 0.3);
    \foreach \x in {-2.0,-1.5,...,2.0} { \draw[gray!50,thin] (\x,-1.2) -- (\x,0.6); }
    \draw[gray!50,thin] (-2.0,0.3) .. controls (-1.5,0.2) and (-0.5,0.0) .. (0,-0.4) .. controls (0.5,0.0) and (1.5,0.2) .. (2.0,0.3);
    \draw[gray!50,thin] (-2.0,0.1) .. controls (-1.5,-0.1) and (-0.5,-0.2) .. (0,-0.6) .. controls (0.5,-0.2) and (1.5,-0.1) .. (2.0,0.1);
    \draw[gray!50,thin] (-2.0,-0.3) .. controls (-1.5,-0.4) and (-0.5,-0.5) .. (0,-0.7) .. controls (0.5,-0.5) and (1.5,-0.4) .. (2.0,-0.3);
    \fill[red!70] (0,-0.55) circle (0.15);
    \node[red] at (0,-0.85) {$M$};
    \draw[yellow!80!orange,very thick] (-2.3,0.8) .. controls (-0.8,0.4) and (-0.4,-0.1) .. (0,-0.3) .. controls (0.4,-0.1) and (0.8,0.4) .. (2.3,0.8);
    \draw[->,yellow!80!orange,very thick] (2.3,0.8) -- (2.7,0.9);
    \draw[->,teal,thick] (0.8,0.5) arc (30:10:2.0);
    \node[teal] at (1.3,0.55) {$\delta\theta$};
    \node at (0,1.0) {光线在引力场中的偏折};
\end{tikzpicture}
\caption{质量 $M$ 使时空弯曲，经过的光子沿弯曲时空中的类光测地线运动。偏折角 $\delta\theta = 4GM/(bc^2)$（$b$ 是瞄准距离），是牛顿预言值的两倍——1919 年日食观测证实了爱因斯坦的预言。}
\label{fig:light_deflection}
\end{figure}

\paragraph{算例 1.2：估算地球表面的时空曲率}

地球表面引力加速度 $g = 9.8\,\text{m/s}^2$，半径 $R_\oplus = 6.4 \times 10^6\,\text{m}$。由 $\Gamma^i{}_{00} \approx \pd_i \Phi$ 和 $\Phi = -GM/r \approx gz$，估得：
\begin{equation}
    |\Gamma| \sim \frac{g}{c^2} \approx \frac{9.8}{(3\times10^8)^2} \approx 10^{-16}\,\text{m}^{-1}.
\end{equation}
曲率张量 $R \sim \pd\Gamma \sim g/(c^2 R_\oplus) \approx 10^{-23}\,\text{m}^{-2}$。地球表面的时空曲率极其微小——只有在强引力场或极高精度测量中才显著。

\paragraph{算例 1.3：GPS 卫星的引力时间膨胀——广义相对论在日常技术中的应用}
\label{para:gps}

作为第 \ref{sec:equivalence_principle} 节中引力时间膨胀推导的直接应用，我们来计算 GPS 卫星与地面时钟之间的固有时差异。这是一个广义相对论效应在日常技术中不可或缺的绝佳例证。

\medskip

\noindent\textbf{已知参数：}
\begin{itemize}
    \item 地球质量：$M_\oplus = 5.97 \times 10^{24}\,\text{kg}$
    \item 地球半径：$R_\oplus = 6.37 \times 10^6\,\text{m}$
    \item GPS 卫星轨道高度：$h = 2.02 \times 10^7\,\text{m}$（距地心距离 $r_{\text{sat}} = R_\oplus + h = 2.66 \times 10^7\,\text{m}$）
    \item 引力常数：$G = 6.67 \times 10^{-11}\,\text{m}^3\text{kg}^{-1}\text{s}^{-2}$
    \item 光速：$c = 3.00 \times 10^8\,\text{m/s}$
\end{itemize}

\medskip

\noindent\textbf{步骤一：计算引力势差。} 在弱场近似下，地球外部的牛顿引力势为：
\begin{equation}
    \Phi(r) = -\frac{GM_\oplus}{r}.
\end{equation}
地面（$r = R_\oplus$）与卫星（$r = r_{\text{sat}}$）之间的引力势差为：
\begin{equation}
    \Delta \Phi = \Phi(r_{\text{sat}}) - \Phi(R_\oplus) = GM_\oplus\left(\frac{1}{R_\oplus} - \frac{1}{r_{\text{sat}}}\right).
\end{equation}
代入数值：
\begin{align}
    \frac{GM_\oplus}{c^2 R_\oplus} &= \frac{(6.67\times10^{-11})(5.97\times10^{24})}{(3.00\times10^8)^2(6.37\times10^6)} \approx 6.95 \times 10^{-10}, \\
    \frac{GM_\oplus}{c^2 r_{\text{sat}}} &= \frac{(6.67\times10^{-11})(5.97\times10^{24})}{(3.00\times10^8)^2(2.66\times10^7)} \approx 1.67 \times 10^{-10}.
\end{align}

\medskip

\noindent\textbf{步骤二：计算引力时间膨胀。} 利用第 \ref{sec:equivalence_principle} 节推导的公式：
\begin{equation}
    \frac{\dd \tau_{\text{sat}}}{\dd \tau_{\text{ground}}} \approx 1 + \frac{\Delta \Phi}{c^2} = 1 + \frac{GM_\oplus}{c^2}\left(\frac{1}{R_\oplus} - \frac{1}{r_{\text{sat}}}\right).
\end{equation}
代入数值：
\begin{equation}
    \frac{\dd \tau_{\text{sat}}}{\dd \tau_{\text{ground}}} \approx 1 + (6.95 \times 10^{-10} - 1.67 \times 10^{-10}) = 1 + 5.28 \times 10^{-10}.
\end{equation}
因此，卫星上的时钟因\textbf{纯引力效应}每天比地面时钟快：
\begin{equation}
    \Delta t_{\text{grav}} = 5.28 \times 10^{-10} \times 86{,}400\,\text{s} \approx 4.56 \times 10^{-5}\,\text{s} = 45.6\,\mu\text{s}.
\end{equation}

\medskip

\noindent\textbf{步骤三：综合效应（狭义相对论修正）。} GPS 卫星以轨道速度 $v \approx 3.87 \times 10^3\,\text{m/s}$ 运动，这产生一个\textbf{相反的}狭义相对论时间膨胀效应：
\begin{equation}
    \frac{\dd \tau_{\text{sat}}}{\dd \tau_{\text{ground}}} \approx 1 - \frac{v^2}{2c^2} \approx 1 - 8.35 \times 10^{-11}.
\end{equation}
即卫星时钟因速度每天慢约 $7.2\,\mu\text{s}$。

\medskip

\noindent\textbf{净效应：}
\begin{equation}
    \boxed{\Delta t_{\text{net}} = \Delta t_{\text{grav}} + \Delta t_{\text{SR}} \approx 45.6\,\mu\text{s} - 7.2\,\mu\text{s} = 38.4\,\mu\text{s}/\text{day}}.
\end{equation}

\medskip

\noindent\textbf{工程意义。} 如果不对此相对论修正进行处理，GPS 的定位误差将以每天约 $c \times 38.4\,\mu\text{s} \approx 11.5\,\text{km}$ 的速度累积——这将使整个系统在几分钟内彻底失效。为此，GPS 卫星在发射前将其原子钟的频率预先调低约 $4.46 \times 10^{-10}$ 倍，以补偿轨道上的净时间膨胀效应。这是广义相对论从抽象理论转化为日常工程实践的经典案例：\textbf{每次你用手机导航时，你都在直接使用爱因斯坦的广义相对论}。


关于符号约定：除非特别说明，本文采用以下约定：
\begin{itemize}
    \item 度规号差为 $(-,+,+,+)$，即闵可夫斯基度规 $\eta_{\mu\nu} = \operatorname{diag}(-1, +1, +1, +1)$
    \item 希腊字母指标 $\mu,\nu,\lambda,\ldots$ 取值 $\{0,1,2,3\}$（时空指标）
    \item 拉丁字母指标 $i,j,k,\ldots$ 取值 $\{1,2,3\}$（空间指标）
    \item 采用 Einstein 求和约定：重复的上指标和下指标自动求和
    \item 黎曼曲率张量定义为 $R^\rho{}_{\sigma\mu\nu} = \pd_\mu \Gamma^\rho{}_{\nu\sigma} - \pd_\nu \Gamma^\rho{}_{\mu\sigma} + \Gamma^\rho{}_{\mu\lambda}\Gamma^\lambda{}_{\nu\sigma} - \Gamma^\rho{}_{\nu\lambda}\Gamma^\lambda{}_{\mu\sigma}$
    \item 速度单位 $c=1$（自然单位制），除非特别注明
\end{itemize}
